\documentclass[10pt,a4paper]{amsart}
\usepackage{amsmath,amssymb,amsthm,graphicx}
% Retained minimal preamble for the Stacks Project chapter more-algebra.tex (More on Algebra).
% The chapter file pinned-original-source.tex carries only \input{preamble} on its first line, so no
% preamble text lives inside the chapter itself.  The Stacks Project uses one single shared preamble for
% every chapter; the snapshot stages it identically as sources/stacks-chow/preamble.tex and
% sources/stacks-etale-cohomology/preamble.tex.  The block below is copied mechanically, line for line,
% from sources/stacks-chow/preamble.tex lines 14-18 (the Xy-pic setup) and 154-217 (the native theorem
% declarations together with the mathematical macros).  Nothing was added, removed or redefined; the
% wrapper class and the amsmath/amssymb/amsthm/graphicx base packages are supplied by the assembler.
% Chapter: more-algebra.tex -- More on Algebra.
% For commutative diagrams we use Xy-pic
\usepackage[all]{xy}

% We use 2cell for 2-commutative diagrams.
\xyoption{2cell}

% Theorem environments.
%
\theoremstyle{plain}
\newtheorem{theorem}[subsection]{Theorem}
\newtheorem{proposition}[subsection]{Proposition}
\newtheorem{lemma}[subsection]{Lemma}

\theoremstyle{definition}
\newtheorem{definition}[subsection]{Definition}
\newtheorem{example}[subsection]{Example}
\newtheorem{exercise}[subsection]{Exercise}
\newtheorem{situation}[subsection]{Situation}

\theoremstyle{remark}
\newtheorem{remark}[subsection]{Remark}
\newtheorem{remarks}[subsection]{Remarks}

\numberwithin{equation}{subsection}

% Macros
%
\def\lim{\mathop{\mathrm{lim}}\nolimits}
\def\colim{\mathop{\mathrm{colim}}\nolimits}
\def\Spec{\mathop{\mathrm{Spec}}}
\def\Hom{\mathop{\mathrm{Hom}}\nolimits}
\def\Ext{\mathop{\mathrm{Ext}}\nolimits}
\def\SheafHom{\mathop{\mathcal{H}\!\mathit{om}}\nolimits}
\def\SheafExt{\mathop{\mathcal{E}\!\mathit{xt}}\nolimits}
\def\Sch{\mathit{Sch}}
\def\Mor{\mathop{\mathrm{Mor}}\nolimits}
\def\Ob{\mathop{\mathrm{Ob}}\nolimits}
\def\Sh{\mathop{\mathit{Sh}}\nolimits}
\def\NL{\mathop{N\!L}\nolimits}
\def\CH{\mathop{\mathrm{CH}}\nolimits}
\def\proetale{{pro\text{-}\acute{e}tale}}
\def\etale{{\acute{e}tale}}
\def\QCoh{\mathit{QCoh}}
\def\Ker{\mathop{\mathrm{Ker}}}
\def\Im{\mathop{\mathrm{Im}}}
\def\Coker{\mathop{\mathrm{Coker}}}
\def\Coim{\mathop{\mathrm{Coim}}}

% Boxtimes
%
\DeclareMathSymbol{\boxtimes}{\mathbin}{AMSa}{"02}

%
% Macros for moduli stacks/spaces
%
\def\QCohstack{\mathcal{QC}\!\mathit{oh}}
\def\Cohstack{\mathcal{C}\!\mathit{oh}}
\def\Spacesstack{\mathcal{S}\!\mathit{paces}}
\def\Quotfunctor{\mathrm{Quot}}
\def\Hilbfunctor{\mathrm{Hilb}}
\def\Curvesstack{\mathcal{C}\!\mathit{urves}}
\def\Polarizedstack{\mathcal{P}\!\mathit{olarized}}
\def\Complexesstack{\mathcal{C}\!\mathit{omplexes}}
% \Pic is the operator that assigns to X its picard group, usage \Pic(X)
% \Picardstack_{X/B} denotes the Picard stack of X over B
% \Picardfunctor_{X/B} denotes the Picard functor of X over B
\def\Pic{\mathop{\mathrm{Pic}}\nolimits}
\def\Picardstack{\mathcal{P}\!\mathit{ic}}
\def\Picardfunctor{\mathrm{Pic}}
\def\Deformationcategory{\mathcal{D}\!\mathit{ef}}

% --- appended 2026-10-12 (Chief): the project own definitions of the reference/slogan/history
% display environments, copied mechanically from the shared Stacks preamble lines 8-12
% (\usepackage{verbatim} plus the three \newenvironment lines, which the chapter preamble keeps
% verbatim from the project preamble). Nothing added, removed or redefined: these environments
% carry no mathematical content and the project itself typesets them as comments. Without these
% three lines a counted statement whose own text uses them fails with Environment slogan undefined.
\usepackage{verbatim}
\newenvironment{reference}{\comment}{\endcomment}
\newenvironment{slogan}{\comment}{\endcomment}
\newenvironment{history}{\comment}{\endcomment}
\usepackage{cleveref}
\crefname{originallemma}{original lemma}{original lemmas}
\Crefname{originallemma}{Original lemma}{Original lemmas}
\makeatletter
\newlabel{lemma-dumb-properties-henselization}{{07QM}{source}{}{sourcepointer-1}{}}
\newlabel{lemma-dumb-properties-henselization@cref}{{[originallemma][0][]07QM}{[1][0][]source}{}{}{}}
\makeatother
\crefname{originallemma}{original lemma}{original lemmas}
\Crefname{originallemma}{Original lemma}{Original lemmas}
\makeatletter
\newlabel{lemma-unibranch}{{0BQ0}{source}{}{sourcepointer-2}{}}
\newlabel{lemma-unibranch@cref}{{[originallemma][0][]0BQ0}{[1][0][]source}{}{}{}}
\makeatother
\crefname{originallemma}{original lemma}{original lemmas}
\Crefname{originallemma}{Original lemma}{Original lemmas}
\makeatletter
\newlabel{lemma-geometrically-unibranch}{{06DM}{source}{}{sourcepointer-3}{}}
\newlabel{lemma-geometrically-unibranch@cref}{{[originallemma][0][]06DM}{[1][0][]source}{}{}{}}
\makeatother
\crefname{originallemma}{original lemma}{original lemmas}
\Crefname{originallemma}{Original lemma}{Original lemmas}
\makeatletter
\newlabel{algebra-lemma-smooth-syntomic}{{00TA}{source}{}{sourcepointer-4}{}}
\newlabel{algebra-lemma-smooth-syntomic@cref}{{[originallemma][0][]00TA}{[1][0][]source}{}{}{}}
\makeatother
\crefname{originallemma}{original lemma}{original lemmas}
\Crefname{originallemma}{Original lemma}{Original lemmas}
\makeatletter
\newlabel{algebra-lemma-flat-localization}{{00HT}{source}{}{sourcepointer-5}{}}
\newlabel{algebra-lemma-flat-localization@cref}{{[originallemma][0][]00HT}{[1][0][]source}{}{}{}}
\makeatother
\crefname{originallemma}{original lemma}{original lemmas}
\Crefname{originallemma}{Original lemma}{Original lemmas}
\makeatletter
\newlabel{algebra-lemma-composition-flat}{{00HC}{source}{}{sourcepointer-6}{}}
\newlabel{algebra-lemma-composition-flat@cref}{{[originallemma][0][]00HC}{[1][0][]source}{}{}{}}
\makeatother
\crefname{originallemma}{original lemma}{original lemmas}
\Crefname{originallemma}{Original lemma}{Original lemmas}
\makeatletter
\newlabel{algebra-lemma-flat-base-change}{{00HI}{source}{}{sourcepointer-7}{}}
\newlabel{algebra-lemma-flat-base-change@cref}{{[originallemma][0][]00HI}{[1][0][]source}{}{}{}}
\makeatother
\crefname{originallemma}{original lemma}{original lemmas}
\Crefname{originallemma}{Original lemma}{Original lemmas}
\makeatletter
\newlabel{algebra-lemma-local-flat-ff}{{00HR}{source}{}{sourcepointer-8}{}}
\newlabel{algebra-lemma-local-flat-ff@cref}{{[originallemma][0][]00HR}{[1][0][]source}{}{}{}}
\makeatother
\crefname{originallemma}{original lemma}{original lemmas}
\Crefname{originallemma}{Original lemma}{Original lemmas}
\makeatletter
\newlabel{algebra-lemma-flat-going-down}{{00HS}{source}{}{sourcepointer-9}{}}
\newlabel{algebra-lemma-flat-going-down@cref}{{[originallemma][0][]00HS}{[1][0][]source}{}{}{}}
\makeatother
\crefname{originallemma}{original lemma}{original lemmas}
\Crefname{originallemma}{Original lemma}{Original lemmas}
\makeatletter
\newlabel{algebra-lemma-henselian-functorial-improve}{{08HU}{source}{}{sourcepointer-10}{}}
\newlabel{algebra-lemma-henselian-functorial-improve@cref}{{[originallemma][0][]08HU}{[1][0][]source}{}{}{}}
\makeatother
\crefname{originallemma}{original lemma}{original lemmas}
\Crefname{originallemma}{Original lemma}{Original lemmas}
\makeatletter
\newlabel{algebra-lemma-strictly-henselian-functorial-improve}{{08HV}{source}{}{sourcepointer-11}{}}
\newlabel{algebra-lemma-strictly-henselian-functorial-improve@cref}{{[originallemma][0][]08HV}{[1][0][]source}{}{}{}}
\makeatother
\crefname{originallemma}{original lemma}{original lemmas}
\Crefname{originallemma}{Original lemma}{Original lemmas}
\makeatletter
\newlabel{algebra-lemma-quotient-henselization}{{05WQ}{source}{}{sourcepointer-12}{}}
\newlabel{algebra-lemma-quotient-henselization@cref}{{[originallemma][0][]05WQ}{[1][0][]source}{}{}{}}
\makeatother
\crefname{originallemma}{original lemma}{original lemmas}
\Crefname{originallemma}{Original lemma}{Original lemmas}
\makeatletter
\newlabel{algebra-lemma-normal-goes-up}{{033C}{source}{}{sourcepointer-13}{}}
\newlabel{algebra-lemma-normal-goes-up@cref}{{[originallemma][0][]033C}{[1][0][]source}{}{}{}}
\makeatother
\crefname{originallemma}{original lemma}{original lemmas}
\Crefname{originallemma}{Original lemma}{Original lemmas}
\makeatletter
\newlabel{algebra-lemma-colimit-normal-ring}{{037D}{source}{}{sourcepointer-14}{}}
\newlabel{algebra-lemma-colimit-normal-ring@cref}{{[originallemma][0][]037D}{[1][0][]source}{}{}{}}
\makeatother
\crefname{originallemma}{original lemma}{original lemmas}
\Crefname{originallemma}{Original lemma}{Original lemmas}
\makeatletter
\newlabel{algebra-lemma-local-tensor-with-integral}{{092Y}{source}{}{sourcepointer-15}{}}
\newlabel{algebra-lemma-local-tensor-with-integral@cref}{{[originallemma][0][]092Y}{[1][0][]source}{}{}{}}
\makeatother
\crefname{originallemma}{original lemma}{original lemmas}
\Crefname{originallemma}{Original lemma}{Original lemmas}
\makeatletter
\newlabel{algebra-lemma-quotient-strict-henselization}{{05WS}{source}{}{sourcepointer-16}{}}
\newlabel{algebra-lemma-quotient-strict-henselization@cref}{{[originallemma][0][]05WS}{[1][0][]source}{}{}{}}
\makeatother
\title{The number of geometric branches is invariant under a smooth local homomorphism}
\author{The Stacks Project Authors}
\date{stacks-project:more-algebra.tex:local-snapshot-20261012; GNU FDL 1.2 or later}
\begin{document}
\maketitle
\setcounter{subsection}{1}
\noindent\textit{Editorial note.} All delivered mathematical text is the author's own, copied from the pinned Stacks Project chapter file more-algebra.tex; no statement or proof step was rewritten, reordered or removed. The counted span is the author environment and its own proof environment. References inside the span are handed to the assembler as navigation stubs that carry the canonical tag of the cited tag only; whether each of them is load-bearing is recorded on the stub itself.\par\medskip
\noindent\textit{Allowed standard prerequisites (editorial).}\begin{itemize}
\item Smooth ring maps, localizations, compositions and base changes of flat maps are flat, and a flat local homomorphism is faithfully flat. used at proof lines 31607-31614. Supplies the flatness and faithful flatness used to compare minimal primes by going down in both parts of the proof.
\item The henselization and the strict henselization are flat, and the (strict) henselization of a base change at a maximal ideal above is the corresponding one of the tensor product. invoked at proof lines 31615-31621 and again at proof lines 31658-31664. Lets the proof replace the rings by their henselizations and keep the hypothesis of being a localization of a smooth map.
\item Going down for flat ring maps lets one compare the minimal primes of a ring and of a flat extension. invoked at proof lines 31622-31628 and again at proof lines 31665-31670. Shows that the minimal primes of the henselization correspond to those of the ring, so the count can be tested prime by prime.
\item The integral closure of a henselian local domain is a local domain, and tensoring it with the henselization of the extension stays normal and local. invoked at proof lines 31633-31652 and again at proof lines 31671-31684. Produces the local normal ring containing the henselization of the larger ring, which forces that henselization to be a domain.
\end{itemize}\medskip
\ifdefined\hypertarget\hypertarget{sourcepointer-1}{}\fi
\noindent\textit{Source reference.} Original lemma 07QM: sources/stacks-more-algebra/pinned-original-source.tex, source lines 10913--10913. Cited in the opening list at proof lines 31613-31614 for flatness of the henselization; the comparison of minimal primes built on that flatness is carried out inside the counted span.\par\medskip
\ifdefined\hypertarget\hypertarget{sourcepointer-2}{}\fi
\noindent\textit{Source reference.} Original lemma 0BQ0: sources/stacks-more-algebra/pinned-original-source.tex, source lines 31249--31249. Invoked at proof lines 31633-31635 for the locality of the integral closure in part (2); the reduction to a domain that precedes it and the tensor product argument that follows are inside the counted span.\par\medskip
\ifdefined\hypertarget\hypertarget{sourcepointer-3}{}\fi
\noindent\textit{Source reference.} Original lemma 06DM: sources/stacks-more-algebra/pinned-original-source.tex, source lines 31453--31453. The analogue of the previous reference for part (1), used at proof lines 31671-31674; the parallel argument in which it is applied is inside the counted span.\par\medskip
\ifdefined\hypertarget\hypertarget{sourcepointer-4}{}\fi
\noindent\textit{Source reference.} Original lemma 00TA: sources/stacks-chow/tags.txt, source lines 1--1. One entry of the opening list at proof line 31608 supplying flatness of the smooth map used throughout; it is a pointer, and every deduction drawn from flatness is inside the counted span.\par\medskip
\ifdefined\hypertarget\hypertarget{sourcepointer-5}{}\fi
\noindent\textit{Source reference.} Original lemma 00HT: sources/stacks-chow/tags.txt, source lines 1--1. One entry of the opening list at proof lines 31609-31610 giving flatness of the localization that the hypothesis mentions; the argument using it is inside the counted span.\par\medskip
\ifdefined\hypertarget\hypertarget{sourcepointer-6}{}\fi
\noindent\textit{Source reference.} Original lemma 00HC: sources/stacks-chow/tags.txt, source lines 1--1. One entry of the opening list at proof lines 31611-31612, used to see that the composite into the henselization is flat; the comparison of minimal primes that uses it is inside the counted span.\par\medskip
\ifdefined\hypertarget\hypertarget{sourcepointer-7}{}\fi
\noindent\textit{Source reference.} Original lemma 00HI: sources/stacks-chow/tags.txt, source lines 1--1. One entry of the opening list at proof lines 31611-31612; it guarantees that the tensor products appearing later are flat, and every such use is inside the counted span.\par\medskip
\ifdefined\hypertarget\hypertarget{sourcepointer-8}{}\fi
\noindent\textit{Source reference.} Original lemma 00HR: sources/stacks-chow/tags.txt, source lines 1--1. One entry of the opening list at proof lines 31612-31613, used to see that every minimal prime below is hit; the prime by prime comparison is inside the counted span.\par\medskip
\ifdefined\hypertarget\hypertarget{sourcepointer-9}{}\fi
\noindent\textit{Source reference.} Original lemma 00HS: sources/stacks-chow/tags.txt, source lines 1--1. Cited twice, in the lists at proof lines 31613-31614 and in both parts at proof lines 31622-31628 and 31665-31670, to compare minimal primes; the comparisons themselves are inside the counted span.\par\medskip
\ifdefined\hypertarget\hypertarget{sourcepointer-10}{}\fi
\noindent\textit{Source reference.} Original lemma 08HU: sources/stacks-chow/tags.txt, source lines 1--1. Used at proof lines 31616-31619 to identify the henselization of the larger ring with that of the tensor product, which justifies replacing the rings by their henselizations; the replacement and what follows are inside the counted span.\par\medskip
\ifdefined\hypertarget\hypertarget{sourcepointer-11}{}\fi
\noindent\textit{Source reference.} Original lemma 08HV: sources/stacks-chow/tags.txt, source lines 1--1. The analogue used at proof lines 31659-31663 for part (1); the parallel argument that follows is inside the counted span.\par\medskip
\ifdefined\hypertarget\hypertarget{sourcepointer-12}{}\fi
\noindent\textit{Source reference.} Original lemma 05WQ: sources/stacks-chow/tags.txt, source lines 1--1. Used at proof lines 31632-31633 to keep the henselian property after dividing out a minimal prime; the reduction to a domain and the subsequent tensor product argument are inside the counted span.\par\medskip
\ifdefined\hypertarget\hypertarget{sourcepointer-13}{}\fi
\noindent\textit{Source reference.} Original lemma 033C: sources/stacks-chow/tags.txt, source lines 1--1. Used at proof lines 31637-31642 to see that the tensor product of the integral closure with the henselization is normal; the deduction that it is a local normal domain is inside the counted span.\par\medskip
\ifdefined\hypertarget\hypertarget{sourcepointer-14}{}\fi
\noindent\textit{Source reference.} Original lemma 037D: sources/stacks-chow/tags.txt, source lines 1--1. Cited parenthetically at proof lines 31641-31643 to extend the normality statement from a localization of an etale algebra to the henselization; the use of that normality is inside the counted span.\par\medskip
\ifdefined\hypertarget\hypertarget{sourcepointer-15}{}\fi
\noindent\textit{Source reference.} Original lemma 092Y: sources/stacks-chow/tags.txt, source lines 1--1. Used at proof lines 31644-31647, where the hypothesis on the residue fields enters, to see that the tensor product is local; the conclusion that it is a local domain and that the henselization is a domain is inside the counted span.\par\medskip
\ifdefined\hypertarget\hypertarget{sourcepointer-16}{}\fi
\noindent\textit{Source reference.} Original lemma 05WS: sources/stacks-chow/tags.txt, source lines 1--1. The analogue of the quotient statement used in part (1) at proof lines 31670-31671; the parallel deduction is inside the counted span.\par\medskip
\begin{lemma}
\label{lemma-invariance-number-branches-smooth}
Let $A \to B$ be a local homomorphism of local rings which
is the localization of a smooth ring map.
\begin{enumerate}
\item The number of geometric branches of $A$ is equal to the number
of geometric branches of $B$.
\item If $A \to B$ induces a purely inseparable extension of residue fields,
then the number of branches of $A$ is the number of branches of $B$.
\end{enumerate}
\end{lemma}

\begin{proof}
We will use that smooth ring maps are flat
(Algebra, Lemma \ref{algebra-lemma-smooth-syntomic}),
that localizations are flat (Algebra, Lemma
\ref{algebra-lemma-flat-localization}),
that compositions of flat ring maps are flat
(Algebra, Lemma \ref{algebra-lemma-composition-flat}),
that base change of a flat ring map is flat
(Algebra, Lemma \ref{algebra-lemma-flat-base-change}),
that flat local homomorphisms are faithfully flat
(Algebra, Lemma \ref{algebra-lemma-local-flat-ff}),
that (strict) henselization is flat
(Lemma \ref{lemma-dumb-properties-henselization}), and
Going down for flat ring maps
(Algebra, Lemma \ref{algebra-lemma-flat-going-down}).

\medskip\noindent
Proof of (2). Let $A^h$, $B^h$ be the henselizations of $A$, $B$. Then $B^h$
is the henselization of $A^h \otimes_A B$ at the unique
maximal ideal lying over $\mathfrak m_B$, see
Algebra, Lemma \ref{algebra-lemma-henselian-functorial-improve}.
Thus we may and do assume $A$ is henselian.
Since $A \to B \to B^h$ is flat, every minimal prime of $B^h$
lies over a minimal prime of $A$ and since $A \to B^h$ is faithfully flat,
every minimal prime of $A$ does lie under a minimal prime of $B^h$;
in both cases use going down for flat ring maps.
Therefore it suffices to show that given a minimal prime
$\mathfrak p \subset A$, there is at most one minimal prime
of $B^h$ lying over $\mathfrak p$.
After replacing $A$ by $A/\mathfrak p$ and $B$ by $B/\mathfrak p B$
we may assume that $A$ is a domain; the $A$ is still henselian by
Algebra, Lemma \ref{algebra-lemma-quotient-henselization}.
By Lemma \ref{lemma-unibranch} we see that the integral closure
$A'$ of $A$ in its field of fractions is a local domain.
Of course $A'$ is a normal domain. By
Algebra, Lemma \ref{algebra-lemma-normal-goes-up}
we see that $A' \otimes_A B^h$ is a normal ring
(the lemma just gives it for $A' \otimes_A B$, to go up to
$A' \otimes_A B^h$ use that $B^h$ is a colimit of \'etale
$B$-algebras and use Algebra, Lemma \ref{algebra-lemma-colimit-normal-ring}).
By Algebra, Lemma \ref{algebra-lemma-local-tensor-with-integral}
we see that $A' \otimes_A B^h$ is local (this is where we
use the assumption on the residue fields of $A$ and $B$).
Hence $A' \otimes_A B^h$ is a local normal ring, hence a local domain.
Since $B^h \subset A' \otimes_A B^h$ by flatness of $A \to B^h$
we conclude that $B^h$ is a domain as desired.

\medskip\noindent
Proof of (1). Let $A^{sh}$, $B^{sh}$ be strict henselizations
of $A$, $B$. Then $B^{sh}$ is a strict henselization of
$A^h \otimes_A B$ at a maximal ideal lying over $\mathfrak m_B$
and $\mathfrak m_{A^h}$, see
Algebra, Lemma \ref{algebra-lemma-strictly-henselian-functorial-improve}.
Thus we may and do assume $A$ is strictly henselian.
Since $A \to B \to B^{sh}$ is flat, every minimal prime of $B^{sh}$
lies over a minimal prime of $A$ and since $A \to B^{sh}$ is faithfully flat,
every minimal prime of $A$ does lie under a minimal prime of $B^{sh}$;
in both cases use going down for flat ring maps.
Therefore it suffices to show that given a minimal prime
$\mathfrak p \subset A$, there is at most one minimal prime
of $B^{sh}$ lying over $\mathfrak p$.
After replacing $A$ by $A/\mathfrak p$ and $B$ by $B/\mathfrak p B$
we may assume that $A$ is a domain; then $A$ is still strictly henselian by
Algebra, Lemma \ref{algebra-lemma-quotient-strict-henselization}.
By Lemma \ref{lemma-geometrically-unibranch} we see that the integral closure
$A'$ of $A$ in its field of fractions is a local domain whose residue
field is a purely inseparable extension of the residue field of $A$.
Of course $A'$ is a normal domain. By
Algebra, Lemma \ref{algebra-lemma-normal-goes-up}
we see that $A' \otimes_A B^{sh}$ is a normal ring
(the lemma just gives it for $A' \otimes_A B$, to go up to
$A' \otimes_A B^{sh}$ use that $B^{sh}$ is a colimit of \'etale
$B$-algebras and use Algebra, Lemma \ref{algebra-lemma-colimit-normal-ring}).
By Algebra, Lemma \ref{algebra-lemma-local-tensor-with-integral}
we see that $A' \otimes_A B^{sh}$ is local (since $A \subset A'$
induces a purely inseparable residue field extension).
Hence $A' \otimes_A B^{sh}$ is a local normal ring, hence a local domain.
Since $B^{sh} \subset A' \otimes_A B^{sh}$ by flatness of $A \to B^{sh}$
we conclude that $B^{sh}$ is a domain as desired.
\end{proof}

\end{document}
